\section*{Resumo}

\noindent
Esta investigação analisa o desempenho educacional no IDEB EFAI 2023 (Índice de Desenvolvimento da Educação Básica, Anos Iniciais do Ensino Fundamental) das escolas da rede municipal, utilizando aprendizado de máquina não supervisionado para identificar padrões associados ao desempenho escolar em contextos de vulnerabilidade socioeconômica. Integrando microdados do Censo Escolar 2023 e indicadores socioeconômicos do IBGE para os 217 municípios do estado com dados disponíveis, adotamos uma abordagem em dois estágios: segmentação municipal via \textit{K-Means} sobre renda \textit{per capita} e IDHM (Índice de Desenvolvimento Humano Municipal), seguida de \textit{Random Forest Classifier} e testes de Mann-Whitney para identificar, dentro de cada grupo, os fatores que diferenciam municípios de alto e baixo desempenho no IDEB. Os resultados identificaram os fatores mais fortemente associados ao desempenho superior. Essas descobertas evidenciam a utilidade do aprendizado não supervisionado para análises educacionais descritivas e oferecem subsídios concretos para políticas educacionais e de alocação de recursos nos municípios. Esta pesquisa contribui para suprir a sub-representação do Maranhão na literatura de mineração de dados educacionais e propõe uma estrutura analítica replicável para outras regiões em desenvolvimento que busquem orientar estratégias de educação pública a partir de dados administrativos.

\medskip
\noindent\textbf{Palavras-chave:} desempenho educacional; IDEB; IDHM; municípios maranhenses; clusterização; \textit{K-Means}; \textit{Random Forest}; políticas educacionais.